\section{Literature Review}
\subsection{Hallucination in Coding Tasks}
Hallucinations have various definitions. One of them was proposed by \cite{tian2025codehalu} together with
\textbf{CodeHalu} benchmark. Authors define hallucination as "syntactically correct and even semantically plausible, but may not execute as expected or fulfill specified requirements".

Systematic review of hallucination in software engineering tasks by \cite{liu2024exploring}
presents complete taxonomy on various hallucination types in code generation tasks.
Additionally authors presented \textbf{HalluCode} benchmark, containing 5563 Python coding tasks
with reference and hallucinated solutions.
Similarly to \cite{liu2024exploring} in \cite{sahoo2024comprehensive} authors provide survey on hallucination detection and
mitigation across various types of models. Authors also provide taxonomy on hallucination dividing
them on four main categories: \textbf{Contextual Disconnection}, \textbf{Semantic Distortion}, \textbf{Content Hallucination}, and \textbf{Factual Inaccuracy}
Another review \cite{huang2025survey} categorizes hallucinations on two types, factuality hallucination and faithfulness hallucinations. In addition 
paper review mitigation strategies, including \textbf{data filtering, RAG, and targeted optimizations} 

Summing all up we can conclude that hallucination refers to situation, then \textbf{models output
have correct code structure but fails to solve specified problem in efficient way, due to misalignment in internal and external knowledge and poor context quality} 

\subsection{Uncertainty Estimations}
One of the ways to detect hallucinations was proposed by \cite{fadeeva2024fact}. Authors presented
a novel fact checking pipeline \textbf{Claim Conditioned Probability} designed especially for
checking of atomic LLM claims.

Another way to detect hallucination firstly was introduced in \cite{azaria2023internalstatellmknows}. SPALMA
uses internal model representations to estimate models uncertainty score. THis paper demonstrated
ability of the models to store uncertainty information inside.
In other paper \cite{duan2024llms} authors explored approach proposed by \cite{azaria2023internalstatellmknows}. Findings show that
the most insights for uncertainty estimations were acquired from models \textbf{middle layers} (10-20 for LLama models).

Such approach demonstrated high performance, hence became vary popular among other researchers like 
\cite{du2024haloscope} that presented HaloScope that also uses features extracted from internal representations, that later passed to simple classifier
Similarly, \cite{shelmanov2025head} trained special \textbf{Uncertainty Quantification Heads} on hidden states of model
to detect hallucinated samples.

In \cite{chuang2024lookback} authors proposed \textbf{Lookback Lense}, a light weight method, that
uses attention weights for input and generated tokens. Proposed approach
matches performance with more complicated detectors, while using simple linear classifier.
Authors of \cite{zhang2025icr} made one step further by analyzing not just internal representations but
their change among layers, proposing \textbf{ICR Probe} and proving a way to distinguish hallucination cause
based on \textbf{ICR Score} value

Interesting methodology was proposed in \cite{srey2026learning}. Authors assumes that dispersion or variance of internal representations across layers
may reflect internal model uncertainty. Proposed approach avoids large resource consumption and shows effective generalization across different domains

\subsection{Uncertainty Estimations in Coding}
In \cite{sharma2025assessingcorrectnessllmbasedcode} authors propose to use uncertainty estimation as a proxy value
for checking LLMs generated code. In paper authors adapted NLP methods for software engineering tasks
and  added a check semantic equivalence check for ensuring that model understood the task.
In contrast \cite{bouchard2026functional} showed that natural-language semantic entropy does not transfer well to code,
instead proposing \textbf{functional equivalence methods}, that include coding analog of semantic entropy.
Another approach \cite{andriushchenko2026efficient} uses combination of features to train small transformer model
for line-level hallucination detection. It outperforms existing UQ baselines across several code domains.

\subsection{Uncertainty Estimations in Agents}
Only a few paper review role of uncertainty estimates on agentic work. 
One of such paper \cite{zhao-etal-2025-uncertainty} proposes \textbf{SAUP}, framework that propagates
uncertainty through each step of agents reasoning, achieving state-of-art performance.
UALA by \cite{han-etal-2024-towards} tried to integrate uncertainty estimate directly into REACT agent. Generated
response evaluated on model uncertainty, and based on result continues with CoT or uses external tool for accessing information

\subsection{Research Gap}
Reviewed papers demonstrated that it is possible to detect model hallucinations in various ways, however current
mitigation techniques manually work with simple one-shot interactions with users. Only limited number of
papers explore influence of uncertainty estimators on coding agents.
Additionally, most of approaches strictly tied to specific architecture, hence for each
new model fine-tuning of classifiers and features from scratch required. 
This thesis ties to address this gap by exploring effect of various uncertainty scores
on agent behaviour in software engineering tasks.
